\subsection{Quark and gluon jets}
\label{subsec:qgjets}

Every emission rate in Sec.~\ref{subsec:coherence} is proportional to the squared
color charge of the emitter, so the two species of initiating parton differ by the
single ratio
\begin{equation}
	\frac{\bm{Q}_g^2}{\bm{Q}_q^2} = \frac{C_A}{C_F}
	= \frac{2N_c^2}{N_c^2-1} = \frac{9}{4},
	\label{eq:casimirRatio}
\end{equation}
a gluon radiating \(9/4\) as copiously as a quark of the same energy. The
consequences are most directly visible in the number of constituents a jet acquires
and in how far from the jet axis they lie.

The multiplicity follows from the angular-ordered cascade of
Eq.~\ref{eq:coherentSudakov}. Writing \(Y=\tfrac12\ln(\tilde{t}/\Lambda^2)\) and using
the running coupling, the double-logarithmic evolution
\(d^2\langle n\rangle/dY^2 = (4C_A/\beta_0Y)\langle n\rangle\) integrates to
\begin{equation}
	\langle n(Y)\rangle \;\propto\;
	\exp\left[\sqrt{\frac{16C_A}{\beta_0}\,Y}\,\right],
	\qquad
	\beta_0 = 11-\tfrac{2}{3}n_f ,
	\label{eq:multiplicityGrowth}
\end{equation}
slower than any power of the jet energy, as anticipated at the end of
Sec.~\ref{subsec:coherence}. The exponent carries \(C_A\) whatever the initiating
parton, since after the first branching the cascade is gluon-dominated; the species
enters only through the normalization, fixed by that first emission, so
\begin{equation}
	\frac{\langle n\rangle_g}{\langle n\rangle_q}
	\;\xrightarrow[\;Y\to\infty\;]{}\; \frac{C_A}{C_F} = \frac{9}{4} .
	\label{eq:multRatio}
\end{equation}
The approach to Eq.~\ref{eq:multRatio} is slow: energy conservation and
next-to-leading corrections to the cascade reduce the ratio to \(\simeq1.5\) at
presently accessible jet energies.

The overall angular size of a jet scales the same way. Adopt the
Sterman--Weinberg criterion --- a jet holds all but a fraction \(\epsilon\) of its
energy within a cone of half-angle \(\delta\) --- for which the probability that a
parton of color charge \(C_i\) qualifies is, at \(\mathcal{O}(\alpha_s)\),
\begin{equation}
	f_i(\epsilon,\delta) = 1 - \frac{4\alpha_s C_i}{\pi}
	\ln\frac{1}{\epsilon}\,\ln\frac{1}{\delta} ,
	\label{eq:stermanWeinberg}
\end{equation}
the double logarithm being the soft-collinear region of Sec.~\ref{subsec:coherence}
once again. Holding \(f\) and \(\epsilon\) fixed and solving for the angular size,
\begin{equation}
	\delta_i \sim \exp\left[-\frac{\pi(1-f)}{4\alpha_s C_i\ln(1/\epsilon)}\right],
	\qquad
	\delta_g \sim \delta_q^{\,C_F/C_A} ,
	\qquad
	\frac{C_F}{C_A} = \frac{4}{9} .
	\label{eq:swWidth}
\end{equation}
Here the charge sits in the denominator of the exponent, so raising \(\delta_q<1\) to
the power \(4/9\) leaves the gluon jet the wider of the two. Since \(1/\alpha_s\)
grows logarithmically with the hard scale, \(\delta\) falls like a negative power of
it, and less rapidly for gluon jets.

The transverse structure is quantified by the generalized angularities of
Chapter~\ref{ch:angularities},
\begin{equation}
	\lambda^\kappa_\beta = \sum_{c\,\in J} z_c^\kappa
	\left(\frac{\Delta R_c}{R}\right)^{\beta},
	\qquad
	z_c = \frac{p_{\mathrm{T},c}}{\sum_{c^\prime\in J}p_{\mathrm{T},c^\prime}},
	\label{eq:angularityDef}
\end{equation}
of which \(\lambda^0_0\) counts constituents and \(\lambda^1_1\) is the jet width, so
that the multiplicity and the width above are one observable at two values of
\((\kappa,\beta)\). Counting emissions in the region \(z\theta^\beta>\lambda\),
instead of outside the Sterman--Weinberg cone, turns the same argument into a
statement about the distribution of the internal structure rather than the overall
size, and the color charge appears once more as a power,
\begin{equation}
	\Sigma_g(\lambda) = \left[\Sigma_q(\lambda)\right]^{C_A/C_F} ,
	\label{eq:casimirPower}
\end{equation}
\(\Sigma_i(\lambda)\) being the fraction of jets with angularity below \(\lambda\).
The exponent is inverted with respect to Eq.~\ref{eq:swWidth} only because that
equation fixes a probability and solves for an angle, where this one fixes the angle
and returns a probability. Chapter~\ref{ch:angularities} develops the point.

Gluon jets are therefore both more populous
and wider than quark jets of the same energy, and the two effects share a single
origin in Eq.~\ref{eq:casimirRatio}. The separation is statistical rather than
absolute --- the distributions overlap heavily --- but it is large enough that
\(\lambda^\kappa_\beta\) and its groomed variants form the basis of quark--gluon
discrimination, and it is why the quark and gluon fractions of a jet sample must be
controlled before any modification of the angular structure can be attributed to a
medium.
